\documentclass[11pt]{article}
\usepackage{docstyle}
\renewcommand\sheetname{Light \textperiodcentered{} In-class work}

\begin{document}
\sheethead{Light: In-class Work}{Year 9 Science \textperiodcentered{} about 25 minutes in four parts}{Name: \blank[45mm]}
Use a sharp pencil and a ruler for every ray. Put \textbf{arrows} on rays and draw virtual rays as \textbf{dotted} lines.

\setpart{A}
\section*{Part A \quad Shadows and colour \hfill\normalfont\small about 7 minutes}

\Q A small lamp shines on a ball in front of a screen. Draw two rays from the lamp that just touch the top and bottom of the ball and reach the screen. Shade the shadow and label the \tUmbra.
\drawbox{66mm}{\TaskShadow[8mm]}

\Q Complete the table.\par\smallskip
\begin{tabularx}{\textwidth}{@{}l l X@{}}
\toprule
\textbf{Object} & \textbf{Light shone on it} & \textbf{Colour it appears}\\
\midrule
white paper & green light & \\[5mm]
blue car & white light & \\[5mm]
blue car & red light & \\[5mm]
yellow banana & green light & \\[5mm]
black shoe & white light & \\[5mm]
\bottomrule
\end{tabularx}

\Q Complete the colour mixing for light.\par\smallskip
red light + blue light = \blank[40mm] \qquad green light + blue light = \blank[40mm]

\newpage
\setpart{B}
\section*{Part B \quad Reflection \hfill\normalfont\small about 7 minutes}

\Q A ray hits a plane mirror at 30\textdegree{} to the mirror surface. Draw the normal and the reflected ray. Label the \tAoI{} $i$ and the \tAoR{} $r$, and write their sizes.
\drawbox{50mm}{\TaskReflect[9mm]{0}}
\par\smallskip $i =$ \blank[25mm] \qquad $r =$ \blank[25mm]

\Q Draw the image A$'$B$'$C$'$ of triangle ABC in the plane mirror. Then draw two rays from A that reflect off the mirror into the eye.
\drawbox{66mm}{\TaskTriangle[8.6mm]}

\Q Complete the periscope: draw the two mirrors and the path of a ray from the object to the eye.
\drawbox{62mm}{\FigPeriscope[8.4mm]{0}}

\newpage
\setpart{C}
\section*{Part C \quad Refraction \hfill\normalfont\small about 6 minutes}

\Q A ray of light hits a glass block. Draw the normal at both faces, the refracted ray inside the block, and the ray that leaves the block.
\drawbox{60mm}{\FigBlock[9mm]{1}}

\Q A spear fisher looks at a fish. Draw the ray of light from the fish to the eye, and mark where the fish appears to be.
\drawbox{62mm}{\TaskFish[9mm]{Sp}}

\Q White light is dispersed by the prism. Which ray, P or Q, is red? Explain.
\drawbox{48mm}{\TaskPrism[9mm]}
\alines{2}

\newpage
\setpart{D}
\section*{Part D \quad The eye and lenses \hfill\normalfont\small about 5 minutes}

\Q Name the parts of the eye labelled A to G.
\begin{center}\FigEye[11mm]{0}\end{center}
\begin{tabularx}{\textwidth}{@{}X X X X@{}}
A \blank[30mm] & B \blank[30mm] & C \blank[30mm] & D \blank[30mm]\\[6mm]
E \blank[30mm] & F \blank[30mm] & G \blank[30mm] & \\
\end{tabularx}

\Q Parallel rays reach a convex lens. Complete the rays after the lens. Mark the \tFocal{} F and the \tFocalLength.
\drawbox{52mm}{\FigConvex[10mm]{2}}

\Q Why can no image be formed at the blind spot?
\alines{2}

\end{document}
